% ==============================================================================
% Doc. 163 — Time as Mode, Mass as Deformation, Fractality as Inside-View
%             of Decompactified Time
% Johann Pascher, May 2026
% ==============================================================================

\chapter*{Doc.~163 -- Time as Mode, Mass as Deformation, Fractality as Inside-View of Decompactified Time}
\addcontentsline{toc}{chapter}{Doc.~163 -- Time as Mode, Mass as Deformation, Fractality as Inside-View}

\section*{Preliminary note}

Doc.~162 established that $T^4$ is not spatially visualisable and need
not be -- the acceptance of mathematical constructs in physics has
not depended on visualisability for over a century, but on
consistency and testable consequences. That is the negative side:
\emph{what visualisation cannot deliver}.

This document makes the positive side explicit: \emph{what an
admitted inside-view metaphor can deliver}. It shows that the
description of $T^4$ remains coherent when time is read not as a
distinguished continuous variable but as a further mode on the
torus -- and that the fractal structure that appears in FFGFT is,
in this reading, not a property of the underlying geometry but a
phenomenon that arises in the transition to decompactified time.
Within this transition -- and only there -- a particular inside-view
metaphor becomes admissible, which we name explicitly and frame
carefully at the end.

% ==============================================================================
\section{Time as mode, not stage}
% ==============================================================================

On $T^4$ all four coordinates are equally periodic. None of them is
distinguished as ``time''; the division into ``three spatial dimensions
plus one temporal dimension'' is not laid down in the carrier itself,
but arises only from the choice of a descriptive perspective. In the
FFGFT reading, what we phenomenally experience as \emph{flowing time}
is a \emph{mode} in the sense of an oscillation or motion pattern on
the torus -- not a stage on which physics takes place.

This position is consistent with two points already established in
the corpus:

\begin{itemize}[leftmargin=2em,itemsep=4pt,topsep=4pt]
  \item In FFGFT, $\hbar$ is an SI conversion factor between energy
        and time measures, not an ontological primitive
        (Doc.~001a, Doc.~230 §\,Operational equivalence). The status
        of time itself is therefore not primitive either, but
        derived.
  \item The foundational relation $\tilde T \cdot m = 1$ is, on
        $T^4$, a strict reciprocity, not an open time coordinate.
        Reciprocity is a property of a closed structure, not of a
        flow.
\end{itemize}

\textbf{Consequence.} What we perceive as ``continuously flowing
time'' is the phenomenal resolution of a discrete torus mode into a
quasi-continuous variable. This resolution we call \emph{decompactification
of time}. It is not wrong -- it delivers exactly the familiar
predictions of standard physics -- but it is a \emph{descriptive
choice}, not a property of the carrier.

% ==============================================================================
\section{Mass as deformation of the carrier}
% ==============================================================================

In general relativity, mass is described as curvature of the
spacetime manifold. On $T^4$, this idea carries over structurally:
what GR reads as local curvature of an open manifold becomes, on the
closed carrier, a local \emph{deformation} of the torus geometry. The
global topology is preserved -- the four identifications continue to
close the carrier -- but the local metric is altered by mass.

Three remarks:

\begin{itemize}[leftmargin=2em,itemsep=4pt,topsep=4pt]
  \item In the limit of small deformation (local inertial frames),
        the deformed torus is indistinguishable from the Euclidean
        spacetime of GR -- GR is, in this reading, the linear
        approximation in which the global closure is invisible.
  \item GR singularities (black hole, Big Bang) appear structurally
        differently on $T^4$ because the carrier is closed through
        identifications (cf.\ Doc.~230 §\,Operational equivalence,
        the point on singularities).
  \item The deformation is not visualisable -- for exactly the
        reasons developed in Doc.~162. What is described here is the
        \emph{structural relationship} between mass and carrier, not
        a picture.
\end{itemize}

% ==============================================================================
\section{Fractality as an inside-view phenomenon}
% ==============================================================================

FFGFT contains a fractal correction structure (factor $K_{\text{frak}}$,
appearing in $\pi$-gates and in the lepton-mass cascade). Until now
this fractality has been treated in the corpus as a geometric
property of $T^4$. The discussion of time-as-mode shows that a more
precise reading is possible:

\begin{quote}
As long as time is a discrete mode on the torus, there is no genuine
self-iteration. A fractal recursion requires the application of an
operation to its own result -- and that requires a direction in
which ``before'' and ``after'' can be distinguished. On a closed
torus mode this direction does not exist. It arises only when the
mode is decompactified into a directed, quasi-continuous time.
\end{quote}

In this reading, the fractal structure is not an ontological feature
of the underlying geometry but a \emph{phenomenon of the
decompactified time description}. It appears as soon as one enters
the mode ``directed time'' -- that is, as soon as one adopts an
inside view in which ``before'' and ``after'' can be distinguished.

\textbf{Structural parallel to the Born rule.} This position is not
isolated. It is structurally parallel to the Born-rule reading from
Doc.~230 (§\,Operational equivalence): there too, the apparent
randomness of quantum measurement arises only in the transition from
the deterministic $T^4$ dynamics to decompactified time. Both
phenomena -- apparent randomness and fractal inside-geometry -- are
consequences of the \emph{same decompactification}, not independent
findings.

This structural parallel is a non-trivial result: two phenomena that
stand unconnected in standard physics (quantum randomness and
fractal corrections) appear in FFGFT as two manifestations of a
single cause.

% ==============================================================================
\section{Inside-view metaphor: brain folds and light}
% ==============================================================================

\textbf{Preliminary note on admissibility.} The following metaphor is
introduced here \emph{explicitly}, because an openly named metaphor
remains controllable, while a hidden allusion invites suspicion of
esotericism. It is an \emph{inside-view metaphor}, not an ontological
claim about $T^4$. It does not say ``$T^4$ looks like a brain''; it
says something more specific and weaker.

\textbf{The metaphor.} Once time is decompactified, an observer
inside the geometry experiences not the simple structure of the
carrier but a fractally folded experiential geometry -- with folds,
fold-curvatures, and self-similarities on several scales. The
folded structure of a brain -- many-times wound, self-similar,
geometrically rich on a small volume -- is a familiar image for
exactly this kind of inside structure. It is not \emph{what} $T^4$
is, but it is an appropriate visualisation of what the
\emph{inside view of decompactified time} sees.

\textbf{The light example sharpens the point.} \emph{Viewed from
outside} -- in the undistorted description -- light moves in a
straight line. Within the deformed geometry, however, the path
actually traversed is different: it follows the folds and windings,
and for the observer situated inside, \emph{this} winding path is
the real one, not the idealised straight line of the outside
description. The standard reading -- straight light paths in
Euclidean space -- is, in this view, the \emph{outside description},
which holds in the limit of undistorted geometry; the inside view
is the description in which we actually live and think.

\textbf{The clean framing -- what the metaphor delivers and what
not.}

\begin{itemize}[leftmargin=2em,itemsep=2pt,topsep=2pt]
  \item \emph{What it delivers:} It illustrates that a fractally
        complex inside-experiential world is compatible with a
        simple outside carrier structure. It makes plausible why
        our phenomenal world \emph{appears} complex while its
        carrier is simple.
  \item \emph{What it is not:} It is not a claim that the universe
        is ``like a brain'', or that consciousness is ``fractal
        cosmic structure'', or anything similar. Such statements
        would overstep the frame from metaphor to ontology and
        would not be defensible within FFGFT.
  \item \emph{Where it applies:} Only in the transition ``discrete
        mode $\to$ decompactified time''. Outside this transition --
        in the pure $T^4$ description -- there is no fractal
        inside-geometry to describe. The metaphor applies to the
        inside view, not to $T^4$ itself.
\end{itemize}

Within this delimitation the metaphor is methodologically admissible
and didactically valuable. It closes a gap that Doc.~162 left open:
Doc.~162 showed what \emph{no} form of visualisation can deliver
(spatial intuition); this document shows what a \emph{deliberately
delimited} inside-view metaphor can deliver (rendering plausible the
relationship between a simple outside structure and a complex
inside world).

% ==============================================================================
\section{Consequences}
% ==============================================================================

From the four preceding sections, three statements follow that
individually already stand in the corpus, but are brought together
here for the first time as a common structural line:

\begin{enumerate}[leftmargin=2em,itemsep=4pt,topsep=4pt]
  \item \textbf{Time is not distinguished.} It is a mode among four
        equally periodic coordinates on $T^4$. Phenomenal
        time-experience is the inside view of decompactified time.

  \item \textbf{Mass is deformation of the carrier.} What GR
        describes as spacetime curvature is, on $T^4$, a local
        deformation of the closed geometry -- with standard GR as
        the linear, locally inertial approximation.

  \item \textbf{Fractality is an inside-view phenomenon.} The
        fractal structure that FFGFT sees appearing in
        $K_{\text{frak}}$, in the lepton-mass cascade, and in the
        $\pi$-gates is a phenomenon of decompactified time, not a
        property of the carrier itself. Parallel to the Born-rule
        reading from Doc.~230, two otherwise unconnected phenomena
        -- quantum randomness and fractal corrections -- thus
        appear as two manifestations of the same cause.
\end{enumerate}

\textbf{Methodological position.} This document tightens no
empirical prediction and changes no calculation. It is a
\emph{reading clarification}: the structural relationship between
the $T^4$ carrier, decompactified time, and phenomenal
experiential world is made explicit. The practical use is that the
metaphor of brain folds -- and similar inside-view images -- now
have a cleanly framed standing in the corpus and need no longer
operate as unexplained allusions.

% ==============================================================================
\section{Summary}
% ==============================================================================

On $T^4$, time is not a stage but a mode. Mass is not an object on
a spacetime but a deformation of the carrier. Fractality is not a
property of the carrier but a phenomenon arising in the transition
to decompactified time -- structurally parallel to the Born rule
and the apparent randomness of quantum measurement. Within this
transition the brain-folds metaphor is an admissible inside-view
visualisation -- not as a claim about $T^4$, but as an image for
how a simple outside structure becomes a fractally wound inside
world.

The outside description is simple. The inside world in which we
live is wound. The two descriptions are not in contradiction --
they are two sides of the same transition.

